% (c) Nikita Lisitsa, lisyarus@gmail.com, 2026

\documentclass[10pt]{beamer}

\usepackage[T2A]{fontenc}
\usepackage[russian]{babel}
\usepackage{minted}

\usepackage[most]{tcolorbox}
\tcbuselibrary{minted}

\usepackage{graphicx}
\graphicspath{ {./images/} }

\usepackage{adjustbox}

\usepackage{color}
\usepackage{soul}
\usepackage{multicol}
\usepackage{multirow}
\usepackage{hyperref}
\usepackage{tikz}
\usetikzlibrary{arrows.meta,positioning,calc}

\usetheme{metropolis}

\definecolor{red}{rgb}{1,0,0}
\definecolor{green}{rgb}{0,0.5,0}
\definecolor{codebg}{RGB}{29,35,49}
\setminted{fontsize=\footnotesize}

\makeatletter
\newcommand{\slideimage}[1]{
  \begin{figure}
    \begin{adjustbox}{width=\textwidth, totalheight=\textheight-2\baselineskip-2\baselineskip,keepaspectratio}
      \includegraphics{#1}
    \end{adjustbox}
  \end{figure}
}
\makeatother

\title{Язык программирования C}
\subtitle{Лекция 6: стандартная библиотка, const}
\author{Лисица Никита}
\date{2026}

\setbeamertemplate{footline}[frame number]

\usemintedstyle{tango}

\begin{document}

\frame{\titlepage}

\begin{frame}[fragile]
\frametitle{Массивы и функции}
\usemintedstyle{lightbulb}
\begin{tcblisting}{listing engine=minted, minted language=cpp, minted options={fontsize=\scriptsize}, listing only, colback=codebg, colframe=codebg, rounded corners}
int* get_array() {
  int array[100];
  for(int i = 0; i < 100; ++i) {
    array[i] = i;
  }
  return array;
}

int main() {
  int *array = get_array();
  printf("%d", array[0]);
}
\end{tcblisting}
\usemintedstyle{tango}
\begin{itemize}
\item Что в этом коде не так? \pause \mintinline{cpp}|array| -- локальная переменная: она живёт на стеке, и будет ``удалена'' после выхода из функции \mintinline{cpp}|get_array()|
\end{itemize}
\end{frame}

\begin{frame}[fragile]
\frametitle{Массивы и функции}
\usemintedstyle{lightbulb}
\begin{tcblisting}{listing engine=minted, minted language=cpp, minted options={fontsize=\scriptsize}, listing only, colback=codebg, colframe=codebg, rounded corners}
int* get_array() {
  int *array = malloc(100 * sizeof(int));
  for(int i = 0; i < 100; ++i) {
    array[i] = i;
  }
  return array;
}

int main() {
  int *array = get_array();
  printf("%d", array[0]);
  free(array); // Не забываем освободить память!
}
\end{tcblisting}
\pause
\usemintedstyle{tango}
\begin{itemize}
\item \mintinline{cpp}|malloc| выделяет память в \textit{куче}, мы сами контролируем момент её освобождения
\end{itemize}
\end{frame}

\begin{frame}[fragile]
\frametitle{Массивы и функции}
\usemintedstyle{lightbulb}
\begin{tcblisting}{listing engine=minted, minted language=cpp, minted options={fontsize=\scriptsize}, listing only, colback=codebg, colframe=codebg, rounded corners}
int remove_odd(int *src, int size, int *dst) {
  int even_count = 0;
  for (int i = 0; i < size; ++i) {
    if ((src[i] % 2) == 0) ++even_count;
  }
  dst = malloc(even_count * sizeof(int));

  int *out = dst;
  for (int i = 0; i < size; ++i) {
    if ((src[i] % 2) == 0) *out++ = src[i];
  }

  return even_count;
}
\end{tcblisting}
\usemintedstyle{tango}
\begin{itemize}
\item Что в этом коде не так? \pause \mintinline{cpp}|dst| -- аргумент функции, в которую \mintinline{cpp}|result| передаётся \underline{по значению}: мы перезаписали \mintinline{cpp}|dst|, но в \mintinline{cpp}|result| осталось старое значение (\mintinline{cpp}|NULL|)
\end{itemize}
\end{frame}

\begin{frame}[fragile]
\frametitle{Массивы и функции}
\usemintedstyle{lightbulb}
\begin{tcblisting}{listing engine=minted, minted language=cpp, minted options={fontsize=\scriptsize}, listing only, colback=codebg, colframe=codebg, rounded corners}
int remove_odd(int *src, int size, int **dst) {
  int even_count = 0;
  for (int i = 0; i < size; ++i) {
    if ((src[i] % 2) == 0) ++even_count;
  }
  *dst = malloc(even_count * sizeof(int));

  int *out = *dst;
  for (int i = 0; i < size; ++i) {
    if ((src[i] % 2) == 0) *out++ = src[i];
  }

  return even_count;
}
\end{tcblisting}
\end{frame}

\begin{frame}[fragile]
\frametitle{Ключевое слово const: переменные}
\usemintedstyle{lightbulb}
\begin{tcblisting}{listing engine=minted, minted language=cpp, minted options={fontsize=\scriptsize}, listing only, colback=codebg, colframe=codebg, rounded corners}
const int page_size = 4096;
const float pi = 3.1415926535f;
\end{tcblisting}
\pause
\usemintedstyle{tango}
\begin{itemize}
\item \mintinline{cpp}|const|-переменной нельзя присвоить другое значение
\pause
\item Защищает программиста от ошибок
\pause
\item Лучше выражает намерение автора кода
\pause
\item Облегчает оптимизацию программы для компилятора
\end{itemize}
\end{frame}

\begin{frame}[fragile]
\frametitle{Ключевое слово const: аргументы функций}
\usemintedstyle{lightbulb}
\begin{tcblisting}{listing engine=minted, minted language=cpp, minted options={fontsize=\scriptsize}, listing only, colback=codebg, colframe=codebg, rounded corners}
void print(const int value) {
  ...
}
\end{tcblisting}
\pause
\usemintedstyle{tango}
\begin{itemize}
\item \mintinline{cpp}|const| относится к самому аргументу функции, т.е. к локальной переменной (а не к тому значению, которое было в него скопировано)
\pause
\item Как правило, смысла в этом мало -- вызывающего функцию не очень волнует, что произойдёт с \textit{копией} его данных
\end{itemize}
\end{frame}

\begin{frame}[fragile]
\frametitle{Ключевое слово const: указатели}
\begin{itemize}
\item В случае указателей \mintinline{cpp}|const| может быть как сам указатель, так и данные, на которые он указывает
\end{itemize}
\usemintedstyle{lightbulb}
\begin{tcblisting}{listing engine=minted, minted language=cpp, minted options={fontsize=\scriptsize}, listing only, colback=codebg, colframe=codebg, rounded corners}
int i = 10, j = 15;

const int * pi = &i; // Указатель на константу
*pi = 20;            // Ошибка: нельзя изменить константное значение
pi = &j;             // Можно переприсвоить

int * const pi = &i; // Константный сам указатель
*pi = 20;            // Можно изменить значение по указателю
pi = &j;             // Нельзя переприсвоить

// Нельзя переприсвоить, нельзя перезаписать значение по указателю
const int * const pi = &i;
\end{tcblisting}
\end{frame}

\begin{frame}[fragile]
\frametitle{Строковые литералы}
\begin{itemize}
\item Следующие два способа инициализации строки делают радикально разные вещи:
\end{itemize}
\usemintedstyle{lightbulb}
\begin{tcblisting}{listing engine=minted, minted language=cpp, minted options={fontsize=\scriptsize}, listing only, colback=codebg, colframe=codebg, rounded corners}
char * str1 = "Hello!";
char str2[] = "World!";
\end{tcblisting}
\pause
\usemintedstyle{tango}
\begin{itemize}
\item Первый создаёт \underline{указатель} на строку
\pause
\item Вторый создаёт \underline{массив символов} с этой строкой
\end{itemize}
\end{frame}

\begin{frame}[fragile]
\frametitle{Строковые литералы}
\begin{itemize}
\item Строки, полученные первым способом, хранятся вместе с глобальными переменными программы, и их \underline{нельзя менять}:
\end{itemize}
\usemintedstyle{lightbulb}
\begin{tcblisting}{listing engine=minted, minted language=cpp, minted options={fontsize=\scriptsize}, listing only, colback=codebg, colframe=codebg, rounded corners}
char * str = "Hello!";
str[0] = 'A'; // Ошибка выполнения
\end{tcblisting}
\pause
\begin{itemize}
\item Строки, полученные втором способом, хранятся на стеке (как обычные переменные), и их \underline{можно менять}:
\end{itemize}
\usemintedstyle{lightbulb}
\begin{tcblisting}{listing engine=minted, minted language=cpp, minted options={fontsize=\scriptsize}, listing only, colback=codebg, colframe=codebg, rounded corners}
char str[] = "World!";
str[0] = 'A'; // Ошибки нет
\end{tcblisting}
\end{frame}

\begin{frame}[fragile]
\frametitle{Строковые литералы}
\begin{itemize}
\item Чтобы на это не попадаться, лучше объявлять такие указатели как \mintinline{cpp}|const|:
\end{itemize}
\usemintedstyle{lightbulb}
\begin{tcblisting}{listing engine=minted, minted language=cpp, minted options={fontsize=\scriptsize}, listing only, colback=codebg, colframe=codebg, rounded corners}
const char * str = "Hello!";

str[0] = 'A'; // Ошибка компиляции: указатель на const
\end{tcblisting}
\end{frame}

\begin{frame}[fragile]
\frametitle{Ключевое слово const: указатели в аргументах функций}
\usemintedstyle{lightbulb}
\begin{tcblisting}{listing engine=minted, minted language=cpp, minted options={fontsize=\scriptsize}, listing only, colback=codebg, colframe=codebg, rounded corners}
size_t strlen(const char * str) {
  ...
}
\end{tcblisting}
\usemintedstyle{tango}
\begin{itemize}
\item Есть ли в этом смысл? \pause Да: функция явно говорит, что не изменит данные по нашему указателю, что в неё можно передать строковой литерал, и т.п.
\pause
\item В таких случаях лучше всегда указывать \mintinline{cpp}|const|!
\end{itemize}
\end{frame}

\begin{frame}[fragile]
\frametitle{East const vs const west}
\begin{itemize}
\item Ключевое слово \mintinline{cpp}|const| можно писать как слева, так и справа от типа -- для компилятора разницы нет:
\end{itemize}
\usemintedstyle{lightbulb}
\begin{tcblisting}{listing engine=minted, minted language=cpp, minted options={fontsize=\scriptsize}, listing only, colback=codebg, colframe=codebg, rounded corners}
// Одно и то же:
const int x = 10;
int const y = 20;

// Одно и то же:
const int * p = &x;
int const * q = &y;
\end{tcblisting}
\pause
\usemintedstyle{tango}
\begin{itemize}
\item \mintinline{cpp}|const| справа называют \textit{east const} (потому восток тоже справа на большинстве карт)
\item \mintinline{cpp}|const| слева называют \textit{const west} (потому запад слева)
\item Оба варианта распространены, это дело вкуса
\end{itemize}
\end{frame}

\begin{frame}[fragile]
\frametitle{Стандартная библиотека: stdlib.h}
\begin{itemize}
\item \mintinline{cpp}|int atoi(const char *)| -- конвертация строки в десятичное число:
\end{itemize}
\usemintedstyle{lightbulb}
\begin{tcblisting}{listing engine=minted, minted language=cpp, minted options={fontsize=\scriptsize}, listing only, colback=codebg, colframe=codebg, rounded corners}
int x = atoi("123");
\end{tcblisting}
\pause
\usemintedstyle{tango}
\begin{itemize}
\item Есть \mintinline{cpp}|atol| для \mintinline{cpp}|long| и \mintinline{cpp}|atoll| для \mintinline{cpp}|long long|
\pause
\item Возвращает 0, если не смог сконвертировать (конкретную ошибку не выдаёт)
\pause
\item Undefined behavior, если значение не влезает в тип возвращаемого значения
\end{itemize}
\end{frame}

\begin{frame}[fragile]
\frametitle{Стандартная библиотека: stdlib.h}
\begin{itemize}
\item \mintinline{cpp}|long strtol(const char *str, char** str_end, int base)| -- конвертация строки в число в любой системе счисления
\pause
\item Указатель на первый символ, не вошедший в число, записывается в \mintinline{cpp}|str_end|
\end{itemize}
\pause
\usemintedstyle{lightbulb}
\begin{tcblisting}{listing engine=minted, minted language=cpp, minted options={fontsize=\scriptsize}, listing only, colback=codebg, colframe=codebg, rounded corners}
char *str = "25hello";
char *end = NULL;
long x = strtol(str, &end, 10);
// end == str + 2
\end{tcblisting}
\end{frame}

\begin{frame}[fragile]
\frametitle{Стандартная библиотека: stdlib.h}
\begin{itemize}
\item \mintinline{cpp}|int rand()| -- генератор случайных чисел
\pause
\item Разные вызовы вернут разные числа \(\Longrightarrow\) использует внутри какое-то глобальное состояние
\pause
\item \mintinline{cpp}|void srand(unsigned int)| -- проинициализировать состояние генератора
\pause
\item При инициализации через \mintinline{cpp}|srand()| одним и тем же значением гарантируется, что \mintinline{cpp}|rand()| будет выдавать одну и ту же последовательность
\pause
\item Не потокобезопасные функции! Лучше избегать глобального состояния и пользоваться специализированными библиотеками
\end{itemize}
\end{frame}

\begin{frame}[fragile]
\frametitle{Стандартная библиотека: stdlib.h}
\begin{itemize}
\item Быстрая сортировка:
\end{itemize}
\usemintedstyle{lightbulb}
\begin{tcblisting}{listing engine=minted, minted language=cpp, minted options={fontsize=\scriptsize}, listing only, colback=codebg, colframe=codebg, rounded corners}
void qsort (void* ptr, size_t count, size_t size,
  int (*compare)(const void*, const void*));
\end{tcblisting}
\pause
\begin{itemize}
\item Двоичный поиск:
\end{itemize}
\begin{tcblisting}{listing engine=minted, minted language=cpp, minted options={fontsize=\scriptsize}, listing only, colback=codebg, colframe=codebg, rounded corners}
void* bsearch(const void *key, const void *ptr, size_t count,
  size_t size, int (*compare)(const void*, const void*));
\end{tcblisting}
\end{frame}

\begin{frame}[fragile]
\frametitle{Стандартная библиотека: stdlib.h}
\begin{itemize}
\item \mintinline{cpp}|int system(const char *)| -- выполнить консольную команду и вернуть её код возврата
\pause
\item \mintinline{cpp}|void exit(int)| -- завершить программу с кодом возврата
\pause
\item \mintinline{cpp}|void abort()| -- завершить программу аварийно (не flush'ит открытые файлы и т.п.)
\end{itemize}
\end{frame}

\begin{frame}[fragile]
\frametitle{Стандартная библиотека: string.h}
\begin{itemize}
\item \mintinline{cpp}|void* memcpy(void *dst, const void *src, size_t count)| -- копирует \mintinline{cpp}|count| байт из \mintinline{cpp}|src| в \mintinline{cpp}|dst|
\pause
\item Крайне оптимизирована, напр. копирует большими кусками через SSE регистры
\pause
\item \mintinline{cpp}|char* strcpy(char *dest, const char *src)| -- скопировать строку (размер понимается через нуль-терминатор)
\end{itemize}
\end{frame}

\begin{frame}[fragile]
\frametitle{Стандартная библиотека: string.h}
\begin{itemize}
\item \mintinline{cpp}|char* strtok(char *str, const char *delim)| -- разбивает строку на части, разделённые символами из строки \mintinline{cpp}|delim|
\pause
\item Первый вызов задаёт строку для разбиения, последующие вызовы делаются с \mintinline{cpp}|str = NULL|
\pause
\item Каждый вызов выдаёт начало следущей части, в её конец \mintinline{cpp}|strtok| запишет нулевой символ (т.е. он меняет входную строку!)
\pause
\item Запоминает текущую строку и позицию в ней в статической переменной внутри
\end{itemize}
\pause
\usemintedstyle{lightbulb}
\begin{tcblisting}{listing engine=minted, minted language=cpp, minted options={fontsize=\scriptsize}, listing only, colback=codebg, colframe=codebg, rounded corners}
char str[] = "To be, or not to be";
char delim[] = " ,";
char *p = strtok(str, delim);
while (p != NULL) {
  printf("%s\n", p);
  p = strtok(NULL, delim);
}
\end{tcblisting}
\end{frame}

\begin{frame}[fragile]
\frametitle{Стандартная библиотека: time.h}
\begin{itemize}
\item \mintinline{cpp}|time_t| -- тип, хранящий время в секундах
\pause
\item Почти всегда -- целочисленный: число секунд с 1 января 1970 (POSIX time)
\pause
\item \mintinline{cpp}|time_t time(time_t *)| -- возвращает текущее время в секундах
\pause
\item Принимает указатель по историческим причинам, можно туда передавать \mintinline{cpp}|NULL|
\end{itemize}
\pause
\usemintedstyle{lightbulb}
\begin{tcblisting}{listing engine=minted, minted language=cpp, minted options={fontsize=\scriptsize}, listing only, colback=codebg, colframe=codebg, rounded corners}
time_t current_time = time(NULL);
printf("POSIX time: %llu", unsigned long long(current_time));
\end{tcblisting}
\end{frame}

\begin{frame}[fragile]
\frametitle{Стандартная библиотека: time.h}
\begin{itemize}
\item \mintinline{cpp}|clock_t| -- тип, хранящий \textit{clock time} -- время, затраченное процессором на нашу программу
\pause
\item Точность не специфицируется; можно сконвертировать в секунды через константу \mintinline{cpp}|CLOCKS_PER_SEC|
\pause
\item \mintinline{cpp}|clock_t clock()| -- возвращает текущее время в \mintinline{cpp}|clock_t|
\pause
\item Если в программе несколько потоков, вернёт суммарное время, затраченное на всех!
\end{itemize}
\pause
\usemintedstyle{lightbulb}
\begin{tcblisting}{listing engine=minted, minted language=cpp, minted options={fontsize=\scriptsize}, listing only, colback=codebg, colframe=codebg, rounded corners}
clock_t start = clock();
big_slow_function();
clock_t end = clock();
long long duration_in_seconds = (end - start) / CLOCKS_PER_SEC;
\end{tcblisting}
\end{frame}

\begin{frame}[fragile]
\frametitle{Стандартная библиотека: time.h}
\begin{itemize}
\item \mintinline{cpp}|struct timespec| -- структура, хранящая время/длительность как пару из секунд и наносекунд
\pause
\item \mintinline{cpp}|int timespec_get(struct timespec *result, int base)| -- вернуть текущее время в часовом поясе / временной системе \mintinline{cpp}|base|
\end{itemize}
\pause
\usemintedstyle{lightbulb}
\begin{tcblisting}{listing engine=minted, minted language=cpp, minted options={fontsize=\scriptsize}, listing only, colback=codebg, colframe=codebg, rounded corners}
struct timespec current_time;
timespec_get(&current_time, TIME_UTC);

long long seconds     = current_time.tv_sec ;
long long nanoseconds = current_time.tv_nsec;
\end{tcblisting}
\end{frame}

\begin{frame}[fragile]
\frametitle{Стандартная библиотека: assert.h}
\begin{itemize}
\item \mintinline{cpp}|assert(...)| -- макрос, проверяющий переданное условие
\pause
\item В случае невыполнения условия в дебажной сборке печатает сообщение об ошибке и завершает программу
\pause
\item В релизной сборке ничего не делает
\pause
\item Позволяет на этапе разработки отловить ошибки, которые дорого проверять в релизной сборке
\end{itemize}
\pause
\usemintedstyle{lightbulb}
\begin{tcblisting}{listing engine=minted, minted language=cpp, minted options={fontsize=\scriptsize}, listing only, colback=codebg, colframe=codebg, rounded corners}
void print_array(int *array, int size) {
  assert(array != NULL);
  ...
}
\end{tcblisting}
\end{frame}

\begin{frame}[fragile]
\frametitle{Стандартная библиотека: assert.h}
\begin{itemize}
\item Как \mintinline{cpp}|assert()| понимает, какая сейчас сборка?
\pause
\item В релизной сборке компилятор добавляет глобальный макрос \mintinline{cpp}|NDEBUG| (\textit{Not DEBUG})
\pause
\item В дебажной сборке не добавляет :)
\end{itemize}
\end{frame}

\begin{frame}[fragile]
\frametitle{Стандартная библиотека: assert.h}
\begin{itemize}
\item \mintinline{cpp}|assert()| стоит использовать только для ошибок программиста / логики программы, а не для валидации данных или ошибок окружения
\end{itemize}
\pause
\usemintedstyle{lightbulb}
\begin{tcblisting}{listing engine=minted, minted language=cpp, minted options={fontsize=\scriptsize}, listing only, colback=codebg, colframe=codebg, rounded corners}
FILE *f = fopen(...);

// Плохая идея! Ошибка от нас не зависит.
// В релизной сборке такую проверку тоже
// нужно делать!
assert(f != NULL);
\end{tcblisting}
\end{frame}

\begin{frame}[fragile]
\frametitle{Стандартная библиотека: stdint.h}
\begin{center}
  \begin{tabular}{ c c }
    \hline
    \mintinline{cpp}|int8_t| & \mintinline{cpp}|uint8_t| \\
    \mintinline{cpp}|int16_t| & \mintinline{cpp}|uint16_t| \\
    \mintinline{cpp}|int32_t| & \mintinline{cpp}|uint32_t| \\
    \mintinline{cpp}|int64_t| & \mintinline{cpp}|uint64_t| \\
    \hline
    \mintinline{cpp}|intmax_t| & \mintinline{cpp}|uintmax_t| \\
    \mintinline{cpp}|intptr_t| & \mintinline{cpp}|uintptr_t| \\
    \hline
  \end{tabular}
\end{center}
\pause
\begin{itemize}
\item \mintinline{cpp}|int32_t| и т.п. -- гарантированно (без-)знаковый, ровно 32 бита (не скомпилируется, если платформа такой тип не поддерживает)
\pause
\item \mintinline{cpp}|intmax_t| -- максимальный поддерживаемый тип
\pause
\item \mintinline{cpp}|intptr_t| -- тип, гарантированно вмещающий любой указатель
\pause
\item Полезно использовать, если важен размер переменной (для сериализации, передачи на GPU, и т.д.)
\end{itemize}
\end{frame}

\end{document}
